% Chapter 6 — Conclusion & Perspectives

\chapter{Conclusion \& Perspectives}

\label{Chapter6}

This work developed a simplified control-oriented model of a ten-cell PEMFC stack,
calibrated it per cell on the Optimization Experiment, derived and analyzed its
linearized state-space representation, characterized its dynamic current response
to load steps, and proposed a PI-based control architecture for the PEMFC power
system.

The main outcomes are: (i)~a six-parameter semi-empirical voltage model coupled to
lumped anode, membrane, and cathode dynamics that fits all healthy cells to within
12\,mV; (ii)~a diagnosis capability --- degraded cells (0, 1, 2) are revealed both
by voltage collapse and by degenerate, physically meaningless fits;
(iii)~an explicit state-space representation in which $A$ and $B$ follow directly
from the dynamic equations while $C$ and $D$ require Jacobian linearization of the
nonlinear voltage equation, leading to proof that the model is generically
controllable while a voltage-only observer can reconstruct the hydrogen pressure
but not the water states; (iv)~an experimental characterization of the current
response to load steps ($\tau_I\approx 8$--$10$\,s,
$T_{s,2\%}\approx 32$--$40$\,s); and (v)~a control architecture based on a DC/DC
boost converter regulated by a PI controller \cite{derbeli2017}, motivated by the
limited authority of the hydrogen supply, in which the converter output is
regulated to behave either as a current source (overcurrent protection) or as a
voltage source (output-voltage regulation).

Perspectives follow directly. On the control side, the next step is to implement
the proposed boost-converter PI loop on the bench and tune it online following
\cite{derbeli2017}, with the identified current timescale bounding the required
closed-loop bandwidth; the observability analysis further supports a
hydrogen-pressure regulator with a reduced-order voltage-based observer, with PI
control as the literature-standard baseline \cite{zhu2017,derbeli2017} and LQ or
time-delay control as upgrades \cite{pukrushpan2003,kim2010}. On the modeling
side, capturing oxygen transients (removing the constant-$P_{O_2}$ assumption),
liquid-water accumulation, and gas-supply delays \cite{kharrat2025} would extend
validity at higher dynamics. Finally, replacing the degraded cells and repeating
the calibration would allow stack-level closed-loop experiments.
